\section[类域论与函数域]{类域论与函数域（2 月 23 日）}

\srcpage{27}

由于德州冬季风暴，我们有一阵子没上课了，所以先回顾一下之前讲到哪里。我们当时正在讨论阿贝尔对偶（abelian duality），以及它如何把主丛（principal bundle）的拓扑与局部系统（local system）的代数几何联系起来。在讨论过的阿贝尔对偶的例子中，有一个（相对较弱的）类域论（class field theory）形式。我们再过一遍。

设 $F$ 是数域（number field），即 $\Q$ 的有限扩张。类域论把 $F$ 的极大无分歧阿贝尔 Galois 扩张的 Galois 群与 $\cO_F$ 的理想类群等同起来\yizhu{数域 $F$ 的「类群」通常指 $\cO_F$ 的理想类群，即分式理想模主分式理想所得的群。}{A（原文 p.27）}。因此可以等同
\[
  \C[\mathrm{Cl}_F] \cong \C[\mathrm{Loc}_1(\Spec \cO_F)].
\]
这个同构的右端可以解释为 Galois 群的特征标（character）。

这并不总是一个有趣的陈述。例如当 $F=\Q$ 时，两边都消失，尽管 $\Q$ 的 Galois 理论很有意思。为了看到这些数据，我们必须允许分歧（ramification），这留待以后再谈；特别地，我们得弄清楚该用什么来替换类群。几何上，如果把 $\Spec F$ 想象成类似 Riemann 曲面（Riemann surface）的东西，那么类域论讲的是覆盖；分歧允许覆盖分支。若真有 Riemann 曲面，允许分支会使问题有趣得多。

不过今天的主要内容，将集中在函数域（function field）这个更几何的设定上。设 $k$ 是域（暂时不加任何条件），$C$ 是 $k$ 上的光滑射影曲线（smooth projective curve）。画图时取 $k=\C$，$C$ 就是紧 Riemann 曲面。对这样的曲线 $C$，我们附上一个域 $F:=k(C)$，即 $C$ 上的有理函数域。除子（divisor）的阿贝尔群 $\Pic(C)$ 落在短正合序列
\[
  0 \longrightarrow \mathrm{Jac}(C) \longrightarrow \Pic(C) \longrightarrow \Z \longrightarrow 0
\]
中，其中到 $\Z$ 的满射由次数（degree）给出：$\sum_{x\in C} n_x x \mapsto \sum_{x\in C} n_x$，$x\in C(k)$。这个短正合序列分裂：取 $x\in C(k)$，则摩天层（skyscraper sheaf）$\Z\cdot x$ 定义了映射 $\Z\to\Pic(C)$，它是一个截面。

这已经与类域论的情形不同：$\Pic(C)$ 是类群的类比物，但它总是无限的，而类群是有限的。

现在设 $k=\F_q$ 有限。则有映射 $\pi_1^{\et}(C)\twoheadrightarrow \widehat{\Z}:=\varprojlim_n \Z/n$。这个 $\widehat{\Z}$ 就是 $\pi_1^{\et}(\Spec k)$，即 $\Gal^{un}(\bar{k}/k)$，拓扑上由 Frobenius 生成。

我们将用（无分歧）Weil 群 $W^{un}(C)$ 替换 $\pi_1^{\et}(C)$，它是拉回
\[
  \begin{tikzcd}
    W^{un}(C) \arrow[r] \arrow[d] & \Z \arrow[d, hook] \\
    \pi_1^{\et}(C) \arrow[r] & \widehat{\Z}.
  \end{tikzcd}
\]

另一件可做的事是取极大阿贝尔商 $W^{ur,ab}(C)$，它满射到 $\Z$，并且对曲线上每个 $k$-点仍有截面。这看起来更像类群。把 $C$ 想成 Riemann 曲面时，$x\in C(k)$ 处的截面有点像把 $\Spec(k)$ 作为圆圈嵌入 3-流形（即曲面丛）中。

至此，用 Weil 群还是 étale 基本群随你。现在我们也可以改动 $\Pic$。$k$ 有限时，Jacobian $\mathrm{Jac}(C)$ 也是有限阿贝尔群，所以所有的「无限性」都来自到 $\Z$ 的次数映射。于是可以把 $\Pic$ 换成它的射有限完备化（profinite completion），即对一切有限子群取极限；它现在满射到 $\widehat{\Z}$。

\srcpage{28}

\begin{theorem}[无分歧类域论]
  存在映射 $\Pic(C)\to\pi_1^{\et,ab}(C)$，它在射有限完备化上是同构。等价地，存在同构
  \[
    \Pic(C)\cong W^{un,ab}(C),
  \]
  它保持次数同态及其截面。
\end{theorem}

这个同构保持相当多的结构。这个陈述并不总有趣：例如当 $C=\P^1$ 时，两边都是 $\Z$。

类域论告诉我们，$\Pic(C)$ 的特征标同构于（无分歧、阿贝尔）Weil 群的特征标。但正如类域论中一样，取特征标时不需要阿贝尔化：这与 $W^{un}(C)$ 的特征标是同一回事，也与 $\pi_1^{\et}(C)$ 的特征标同一回事——不过这里必须限制在有限阶特征标上。于是映射到 $\Q/\Z$，而不是 $\C^\times$。而 $\pi_1$ 的特征标理应称为 $C$ 上的秩一局部系统（同样，系数上还有些问题；对有有限系数的系统当然成立）。所以 $\Pic^\vee(C)\cong \mathrm{Loc}(C)$。

在 $\Pic(C)$ 上的函数之中，有特征标：它们知道并尊重乘法律；这些被等同于局部系统集合上的函数。但这不是思考它们的唯一方式。

回忆我们可以通过拉回—推前图表来想特征标，就像前一阵讨论过的；这让你看到，它们恰好是 $\Pic(C)$ 在自身上的平移作用的特征函数（eigenfunction）。又因为 $\Pic(C)$ 由 $\Z x$（$x\in C(k)$）生成，特征标可以看成 $\Pic(C)$ 上的函数，它们对每个 $x\in C(k)$ 的 $\Z x$ 作用都是特征函数。这个作用具体是什么？给定 $\cL\in\Pic(C)$ 与 $x\in C$，作用把 $\cL\mapsto\cL(x)$：在 $\cL$ 的除子上加 $1\cdot x$。

「是特征函数」这件事告诉你：函数 $f$ 在这个过程中行为良好：
\[
  f(\cL(x)) = \gamma_x\cdot f(\cL).
\]

现在，通过 Fourier 变换，$\Pic(C)$ 上的函数被等同于 $\mathrm{Loc}(C)$ 上的函数。设 $\rho$ 是局部系统，看成映射 $\rho\colon\pi_1^{\et}(C)\to F$，其中 $F$ 是某个系数域。这给我们一堆本征数据：
\[
  f(\cL(x)) = \rho(\Frob_x) f(\cL).
\]

算子 $\cL\mapsto\cL(x)$ 是 Hecke 算子（Hecke operator）的第一个例子，类似于我们在物理中见到的 't Hooft 算子或 Dirac 单极子（monopole）算子：那里把 3-流形 $M$ 的 $H^2$ 等同于一个格，并通过在格上平移来作用。我们还把 $H^2(M)$ 看成 $\pi_0(\Pic(M))$。

这里 Wilson 算子的类比物是 $\rho\in\mathrm{Loc}(C)\mapsto\rho(\Frob_x)$：在 3-流形 $M$ 上，我们把局部系统送到沿每个环路的单值（monodromy）所给出的乘法。

在代数几何中，我们其实不能把局部系统说成从 $\pi_1$ 出发的映射，所以会避开无限覆盖。有限覆盖对应于映射 $\pi_1\to\mu_n\subset\Q/\Z$。而在特征 $p$ 中，理论只在 $p\nmid n$ 时表现良好。所以我们将取与 $p$ 不同的素数 $\ell$；这时有漂亮的 $\ell$-进局部系统理论。具体地说，考虑映射 $\pi_1\to\Z/\ell^n$，它落在 $\Zl$ 中；我们甚至可以一路走到 $\Q_\ell$ 或 $\Ql$。有很好的 $\pi_1$ 的 $\ell$-进表示理论，例如取值于 $\mathrm{GL}_n(\Ql)$。

$\Ql$ 的好处是存在一个（高度不典范、难以描述的）同构 $\Ql\cong\C$，基本上因为两者都是特征零的代数闭域。而这个同构当然不尊重拓扑！

所以，我们所说的局部系统，指的是从 $C$ 的无分歧 Galois 群到 $\Ql^\times$ 的连续表示（在射有限拓扑下）。在这种情形下，凡是以前说「复值函数」的地方，现在都指 $\Ql$-值函数。故事的形状是一样的，哪怕细节不同。

我们的推广的关键，在于如何认定 $\Pic(C)$。我们一直把它想成除子模主除子所得的群，即商
\[
  \Bigl(\bigoplus_{x\in C}\Z\Bigr)\big/k(C)^\times.
\]

这与线丛的联系如下：给定线丛 $\cL$，可以一般地平凡化它，即在除有限多个点 $x_1,\dots,x_n$ 之外处处有 $\cL$ 的光滑截面，从而定义同构 $\cO_C|_{C\setminus\{x_i\}}\cong\cL|_{C\setminus\{x_i\}}$。另一方面，对每个 $x\in C$，可以在 $x$ 附近很近处平凡化 $\cL$。这里指的是取 $x$ 周围的一个圆盘 $D_x:=\Spec(\cO_x)$（$\cO_x$ 表示 $x$ 处的完备局部环，选一个坐标就把它等同于 $k[[t]]$）。$\cL$ 到 $D_x$ 的拉回被平凡化。

\srcpage{29}

截面 $s$ 在 $x$ 处的局部次数 $\deg_x s$ 非零；但通过改变 $\cL$ 在 $D_x$ 上的平凡化，我们得到「同样的」局部次数。所以我们真正该看的是 $K_x^\times/\cO_x^\times$。选一个坐标就定义了同构 $k((t))^\times/k[[t]]^\times$。具体地说，这是系数非零的 Laurent 级数的等价类，两个级数若在乘以一个可逆 Taylor 级数后一致，就视为等价。这意味着可以把除首项 $t^{-N}$ 之外的所有项都消掉，于是把这个商等同于 $\Z$——次数正该住在这里。很好！

现在我们要变动线丛。这意味着必须变动平凡化所在的轨迹 $U:=C\setminus\{x_i\}$。我们用所有点来处理这件事，这会带入大量数据：一个线丛 $\cL$、一个有理截面，以及对所有 $x\in C$ 的 $\cL|_{D_x}$ 的平凡化。这是对一切 $x\in C$ 的 $K_x^\times$ 的巨大乘积，但我们只看到这个乘积中的一部分元素。这就是限制乘积（restricted product）
\[
  {\prod_{x\in C}}' K_x^\times \subset \prod_{x\in C} K_x^\times,
\]
意思是：除有限多个 $x\in C$ 之外，我们都得到一个元素 $\cO_x^\times\subset K_x^\times$。这仍然是很多数据！

所以现在来缩减数据。

现在，若取线丛与有理截面，最终得到除子，这令人安心：我们得到
\[
  {\prod_{x\in C}}' K_x^\times \Big/ \prod_{x\in C}\cO_x^\times
  = {\prod_{x\in C}}'(K_x^\times/\cO_x^\times)
  = {\prod_{x\in C}}' \Z,
\]
即 $x\in C$ 的有限形式和，这正是除子。（TODO：我需要修正这些限制乘积。）

于是可以把 Picard 群等同于一个双商
\[
  F^\times \Big\backslash {\prod_{x\in C}}' K_x^\times \Big/ \prod_{x\in C}\cO_x^\times,
\]
其中左边的 $F^\times$ 负责截面的改变，右边的乘积负责 $D_x$ 上平凡化的改变。这就是除子模主除子（$F^\times$）所得的群。在数域设定中，这对应于 $F$ 的（无分歧）idele 类群。
